% MT060 / answer-zh — canonical editable content.
(a) \textbf{\(C\)较大}时，违反间隔约束的代价很高。随卷答案左图的边界会弯曲，以适应孤立的红点。这个极限能达到零训练误差，前提是数据在所选二次特征空间中可分。


(b) 对\textbf{严格的\(C\to0^+\)极限}，采用各样本权重相同、偏置不受正则惩罚的标准软间隔目标。原图有10个红点、7个绿点。令红类为\(-1\)、绿类为\(+1\)，候选解\(w=0,b=-1\)的目标值为\(14C\)，因此最优解满足\(\|w_C\|^2\le28C\)，权重趋于0。继续使hinge损失最小，会让\(b_C\to-1\)，即多数类的符号。因此，图示区域最终全部判为红类，区域内不再有分隔边界。随卷右图表示较小的正\(C\)，不是这个严格极限。


(c) 如果这些孤立点是噪声，取适中的正\(C\)可能比努力拟合它们更容易推广到新数据，但不能因此说全判红类的极限最好。仅凭训练图无法确定测试准确率，应在验证数据上选择\(C\)。
